\section{集中不等式：Markov、Chebyshev 与 Chernoff}


这一节把“偏离期望的概率有多小”变成可计算的不等式。三个层次依次是：

1. \textbf{Markov}：只用非负性；
2. \textbf{Chebyshev}：用方差；
3. \textbf{Chernoff}：用矩生成函数/指数变换，得到指数级尾界。

\textbf{本节用到的记号：} $X$ 为随机变量；$\mu=E[X]$；$\Delta=\sqrt{\operatorname{Var}(X)}$；$D(q\Vert p)$ 为二元 KL；$X_1,\dots,X_n$ 独立同分布。

\subsection{Markov 不等式}


若 $X\ge0$，则对任意 $a>0$，

\[
\boxed{
P(X\ge a)\le\frac{E[X]}{a}
}.
\]

等价地，取 $a=\alpha E[X]$：

\[
P(X\ge\alpha E[X])\le\frac1\alpha\qquad(\alpha>0).
\]

\textbf{证明.} 因为 $X\ge a\mathbf 1_{\{X\ge a\}}$，取期望：

\[
E[X]\ge aP(X\ge a).
\]

两边除以 $a$ 即得。

\textbf{用途.} Markov 只要求非负，条件最弱，但界通常很松；它是后面 Chebyshev 和 Chernoff 的证明工具。

\subsection{Chebyshev 不等式}


对任意 $a>0$，

\[
\boxed{
P(|X-E[X]|\ge a)\le\frac{\operatorname{Var}(X)}{a^2}
}.
\]

若记 $\Delta=\sqrt{\operatorname{Var}(X)}$，则

\[
P(|X-E[X]|\ge a\Delta)\le\frac1{a^2}.
\]

\textbf{证明.} 对非负随机变量 $(X-E[X])^2$ 用 Markov：

\[
P((X-E[X])^2\ge a^2)
\le\frac{E[(X-E[X])^2]}{a^2}
=\frac{\operatorname{Var}(X)}{a^2}.
\]

\textbf{用途.} 只要知道方差，就能给出“偏离 $a$ 个标准差”的概率上界；但尾界是多项式级（$1/a^2$），不够精细。

\subsection{Chernoff 不等式}


对任意 $a$ 和任意 $t>0$，

\[
\boxed{
P(X\ge a)
\le\min_{t>0}\frac{E[e^{tX}]}{e^{ta}}
}.
\]

\textbf{证明.} 对非负随机变量 $e^{tX}$ 用 Markov：

\[
P(X\ge a)=P(e^{tX}\ge e^{ta})
\le\frac{E[e^{tX}]}{e^{ta}}.
\]

对所有 $t>0$ 取最小，得到最紧的界。

\textbf{为什么强.} $E[e^{tX}]$ 把整个分布的信息都编码进来；选择最优 $t$ 后，尾界通常是指数级 $e^{-\Theta(n)}$，比 Chebyshev 的 $1/a^2$ 精细得多。

\subsection{Bernoulli 和的 Chernoff 界}


设 $X_1,\dots,X_n$ 独立同分布 $\mathrm{Bern}(p)$，$q>p$。则

\[
P\Bigl(\sum_{i=1}^nX_i\ge qn\Bigr)
\le\exp\bigl(-nD(q\Vert p)\bigr),
\]

其中

\[
D(q\Vert p)
=q\log\frac qp+(1-q)\log\frac{1-q}{1-p}.
\]

\textbf{推导.} 对任意 $t>0$，

\[
\begin{aligned}
P\Bigl(\sum_iX_i\ge qn\Bigr)
&\le\frac{E[e^{t\sum_iX_i}]}{e^{tqn}}\\
&=\frac{\bigl(pe^t+(1-p)\bigr)^n}{e^{tqn}}.
\end{aligned}
\]

令 $f(t)=\log\bigl(pe^t+(1-p)\bigr)-tq$，求导并令其为零：

\[
\frac{pe^t}{pe^t+(1-p)}=q
\quad\Longrightarrow\quad
e^t=\frac{q(1-p)}{p(1-q)}.
\]

代回并整理，指数变为

\[
-\Bigl[q\log\frac qp+(1-q)\log\frac{1-q}{1-p}\Bigr]
=-D(q\Vert p).
\]

\textbf{二次下界.} 采用自然对数时，当 $q,p$ 接近，

\[
D(q\Vert p)\ge 2(p-q)^2.
\]

因此

\[
P\Bigl(\sum_iX_i\ge qn\Bigr)
\le\exp\bigl(-2n(p-q)^2\bigr).
\]

若改用 $\log_2$，$D$ 的数值会放大 $1/\ln2$ 倍，$\exp$ 也要换成以 $2$ 为底；不要混用底数。手写稿里的 $2\log e$ 就是换底因子。

\subsection{与 CLT / 正态尾的比较}


对 $p=1/2$，中心极限定理给出

\[
\frac{\sum_iX_i-n/2}{\sqrt{n}/2}
\approx \mathcal N(0,1).
\]

因此 $P(\sum_iX_i\ge qn)$ 的正态近似尾概率约为

\[
\exp\bigl(-2n(q-1/2)^2\bigr),
\]

与 Chernoff 界的指数一致。

\textbf{但要区分“近似”和“上界”：}

\begin{itemize}
\item CLT 给的是\textbf{近似}，不是严格的尾上界；
\item 严格的指数上界用 Chernoff；
\item 若要一个简单的严格界，也可用 Hoeffding：
\end{itemize}

\[
  P(|\bar X-E[X]|\ge a)\le2e^{-2na^2}.
\]

手写稿把“CLT 正态尾”和“Chernoff 界”放在一起比较，是为了说明两者的指数相同；证明严格上界时不要把 CLT 当成不等式使用。

\subsection{三种界的关系}


\begin{itemize}
\item \textbf{Markov}：需要“非负、期望”；尾界量级 $1/a$；最粗略。
\item \textbf{Chebyshev}：需要“方差”；尾界量级 $1/a^2$；快速估计。
\item \textbf{Chernoff}：需要“矩生成函数”；尾界量级 $e^{-\Theta(n)}$；独立和、大偏差。
\end{itemize}


\textbf{记忆.} Markov 看“质量”，Chebyshev 看“方差”，Chernoff 看“指数矩”；条件越强，界越紧。

\subsection{常见错误}


\begin{itemize}
\item 对非正随机变量直接用 Markov。
\item 把 Chernoff 写成 $P(X\ge a)\le e^{-ta}E[e^{tX}]$ 后忘记对所有 $t>0$ 取最小。
\item 把 CLT 的近似尾概率当成严格上界。
\item 在 Bernoulli 界中把 $D(q\Vert p)$ 的两个参数顺序写反：$D(q\Vert p)$ 是“用 $p$ 近似 $q$”的代价，不是 $D(p\Vert q)$。
\end{itemize}
